% Part: sets-functions-relations
% Chapter: functions
% Section: inverses

\documentclass[../../../include/open-logic-section]{subfiles}

\begin{document}

\olfileid{sfr}{fun}{inv}
\olsection{फलनांचे व्युत्क्रम}

\begin{explain}
फलनांचा विचार आपण संगती म्हणून करतो. मग ही संगती ``उलटवता'' येते का,
हा प्रश्न स्वाभाविकपणे पडतो. उदाहरणार्थ, उत्तरवर्ती फलन $f(x) = x + 1$
उलटवता येते, या अर्थाने की $g(y) = y - 1$ हे फलन $f$ ने केलेली
क्रिया ``पूर्ववत करते''.

पण येथे काळजी घेतली पाहिजे. ~$g$ च्या व्याख्येने $\Int \to \Int$
असे फलन ठरत असले, तरी $\Nat \to \Nat$ असे \emph{फलन} ठरत नाही,
कारण $g(0) \notin \Nat$. म्हणून साध्या उदाहरणांतसुद्धा फलन उलटवता
येते का हे तितकेसे उघड नसते; ते प्रांत आणि सहप्रांतावर अवलंबून असू शकते.

फलनाच्या व्युत्क्रमाच्या संकल्पनेने हा विचार अधिक नेमका करता येतो.
\end{explain}

\begin{defn}
सर्व $x \in A$ आणि $y \in B$ साठी $f(g(y)) = y$ आणि $g(f(x)) = x$
असेल, तर $g \colon B \to A$ हे फलन $f \colon A \to B$ या फलनाचा
\emph{व्युत्क्रम} असते.
\end{defn}

$f$ ला व्युत्क्रम~$g$ असेल, तर आपण अनेकदा~$g$ ऐवजी $f^{-1}$ असे लिहितो.

\begin{explain}
आता कोणत्या फलनांना व्युत्क्रम असतात हे ठरवू. $f\colon A \to B$
च्या व्युत्क्रमासाठी एक संभाव्य फलन म्हणजे पुढीलप्रमाणे ``व्याख्या केलेले''
$g\colon B \to A$:
\[
g(y) = \text{$f(x) = y$ असणारा ``तो'' $x$.}
\]
पण ``व्याख्या केलेले'' आणि ``तो'' यांभोवतीची अवतरणचिन्हे सुचवतात की
ही खरे तर व्याख्या नाही. किमान पूर्ण व्यापकतेने ती नेहमी लागू पडणार
नाही. कारण या व्याख्येने फलन ठरवायचे असेल, तर $f(x) = y$ असणारा
एक आणि केवळ एकच~$x$ असला पाहिजे---~$g$ चे प्रदान एकमेव ठरले पाहिजे.
शिवाय ते प्रत्येक $y \in B$ साठी ठरले पाहिजे. $x_1$ आणि $x_2 \in A$
हे $x_1 \neq x_2$ असूनही $f(x_1) = f(x_2)$ असतील, तर
$y = f(x_1) = f(x_2)$ साठी $g(y)$ एकमेव ठरणार नाही.
आणि $f(x) = y$ असणारा~$x$ मुळीच नसेल, तर $g(y)$ मुळीच ठरणार नाही.
दुसऱ्या शब्दांत, $g$ परिभाषित होण्यासाठी $f$ हे !!{injective} आणि
!!{surjective} दोन्ही असले पाहिजे.

हे टप्प्याटप्प्याने पाहू. प्रश्नाचे दोन भाग करू: फलन~$f\colon A \to B$
दिले असता, $g(f(x)) = x$ असणारे फलन $g\colon B \to A$ कधी असते?
असे $g$ हे $f$ ने केलेली क्रिया ``पूर्ववत करते'', आणि त्याला~$f$ चा
\emph{डावा व्युत्क्रम} म्हणतात. दुसरे म्हणजे, $f(h(y)) = y$ असणारे
फलन $h\colon B \to A$ कधी असते? अशा $h$ ला~$f$ चा
\emph{उजवा व्युत्क्रम} म्हणतात---$f$ हे~$h$ ने केलेली क्रिया ``पूर्ववत करते''.
\end{explain}

\begin{prop}
A अरिक्त असेल आणि $f\colon A \to B$ हे !!{injective} असेल,
तर~$f$ चा एक
\emph{डावा व्युत्क्रम}~$g\colon B \to A$ असतो, आणि सर्व $x \in A$
साठी $g(f(x)) = x$ असते.
\end{prop}
%READERNOTE{OLFUN-001 — गोठवलेल्या इंग्रजी प्रमेयात A अरिक्त असण्याची अट सुटली आहे. रिक्त A पासून अरिक्त B कडेचे रिक्त फलन एकास-एक असूनही B पासून A कडे फलन नसते. नेमकी सामान्य अट: एकास-एक फलनाला डावा व्युत्क्रम तेव्हा आणि केवळ तेव्हाच असतो, जेव्हा A अरिक्त असतो किंवा B रिक्त असतो. मराठी प्रमेयात साधी पुरेशी A अरिक्त ही अट जोडली आहे.}

\begin{proof}
A अरिक्त आहे आणि $f\colon A \to B$ हे !!{injective} आहे असे समजा.
$y \in B$ विचारात घ्या.
$y \in \ran{f}$ असेल, तर $f(x) = y$ असणारा एक $x \in A$ असतो.
$f$ हे !!{injective} असल्यामुळे असा~$x \in A$ एकच असतो. म्हणून
$g(y) = x$ अशी व्याख्या करता येते; म्हणजे $g(y)$ हा $f(x) = y$
असणारा ``तो'' $x \in A$ असतो. $y \notin \ran{f}$ असेल, तर त्याची
संगती कोणत्याही~$a \in A$ शी लावता येते. म्हणून एक $a \in A$
निवडून $g\colon B \to A$ ची व्याख्या अशी करता येते:
\[
g(y) = \begin{cases}
    x & \text{जर $f(x) = y$ असेल}\\
    a & \text{जर $y \notin \ran{f}$ असेल.}
\end{cases}
\]
हे सर्व $y \in B$ साठी परिभाषित आहे, कारण अशा प्रत्येक $y \in \ran{f}$
साठी $f(x) = y$ असणारा नेमका एक $x \in A$ असतो. व्याख्येनुसार,
$y = f(x)$ असल्यास $g(y) = x$, म्हणजे $g(f(x)) = x$.
\end{proof}

\begin{prob}
$f\colon A \to B$ ला डावा व्युत्क्रम~$g$ असल्यास $f$ हे
!!{injective} असते, हे दाखवा.
\end{prob}

\begin{prop}
    $f\colon A \to B$ हे !!{surjective} असेल, तर~$f$ चा एक
    \emph{उजवा व्युत्क्रम}~$h\colon B \to A$ असतो, आणि सर्व~$y \in B$
    साठी $f(h(y)) = y$ असते.
\end{prop}

\begin{proof}
$f\colon A \to B$ हे !!{surjective} आहे असे समजा. $y \in B$ विचारात घ्या.
~$f$ हे !!{surjective} असल्यामुळे $f(x_y) = y$ असणारा एक $x_y \in A$
असतो. म्हणून $h(y) = x_y$ अशी व्याख्या करता येते: प्रत्येक $y \in B$
साठी $f(x) = y$ असणारा काही $x \in A$ आपण निवडतो. ~$f$ हे
!!{surjective} असल्यामुळे निवडण्यासाठी नेहमी किमान एक घटक असतो.\footnote{
$f$ हे !!{surjective} असल्यामुळे प्रत्येक~$y \in B$ साठी
$\Setabs{x}{f(x) = y}$ हा संच अरिक्त असतो. ~$h$ ची आपली व्याख्या
या प्रत्येक संचातून एकच $x$ निवडण्याची मागणी करते. हे नेहमी करता येते
ही गोष्ट खरे तर उघड नाही---या निवडी करता येतात हे एक स्वयंसिद्धक म्हणून
गृहीत धरले जाते. दुसऱ्या शब्दांत, ही प्रतिज्ञा तथाकथित निवडीचे स्वयंसिद्धक
गृहीत धरते. या प्रश्नाकडे आपण \oliflabeldef{sth:choice::chap}{\olref[sth][choice][]{chap} मध्ये पुन्हा वळू}{सध्या दुर्लक्ष करू}.
मात्र अनेक विशिष्ट उदाहरणांत, जसे $A = \Nat$ असेल किंवा सांत असेल,
किंवा $f$ हे !!{bijective} असेल, तेव्हा निवडीचे स्वयंसिद्धक लागत नाही.
($f$ हे !!{bijective} असण्याच्या विशिष्ट उदाहरणात प्रत्येक $y \in B$
साठी $\Setabs{x}{f(x) = y}$ या संचात नेमका एक !!{element} असतो,
म्हणून निवड करण्याचा प्रश्नच उरत नाही.)}
व्याख्येनुसार, $x = h(y)$ असल्यास $f(x) = y$ असते; म्हणजे कोणत्याही
$y \in B$ साठी $f(h(y)) = y$.
\end{proof}

\begin{prob}
$f\colon A \to B$ ला उजवा व्युत्क्रम~$h$ असल्यास $f$ हे
!!{surjective} असते, हे दाखवा.
\end{prob}

\begin{explain}
  मागील सिद्धतेतील कल्पना एकत्र केल्यावर प्रत्येक !!{bijection} ला
  व्युत्क्रम असतो हे मिळते. म्हणजे~$f$ चा डावा आणि उजवा व्युत्क्रम
  असे दोन्ही असणारे एकच फलन असते.
\end{explain}

\begin{prop}\ollabel{prop:bijection-inverse}
$f\colon A \to B$ हे !!{bijective} असेल, तर एक
फलन~$f^{-1}\colon B \to A$ असते, आणि सर्व $x \in A$ साठी
$f^{-1}(f(x)) = x$ तसेच सर्व $y \in B$ साठी $f(f^{-1}(y)) = y$ असते.
\end{prop}

\begin{proof}
सरावासाठी.
\end{proof}

\begin{prob}
\olref[sfr][fun][inv]{prop:bijection-inverse} सिद्ध करा. ~$f^{-1}$ ची
व्याख्या करा, ते फलन आहे हे दाखवा आणि ते~$f$ चा व्युत्क्रम आहे हे
दाखवा. म्हणजे, सर्व $x \in A$ आणि $y \in B$ साठी
$f^{-1}(f(x)) = x$ आणि $f(f^{-1}(y)) = y$ हे दाखवा.
\end{prob}

\begin{explain}
व्युत्क्रम मिळवण्याची थोडी अधिक व्यापक पद्धत आहे. \olref[kin]{sec}
मध्ये आपण पाहिले की प्रत्येक फलन $f$ पासून सर्व $x \in A$ साठी
$f'(x) = f(x)$ असे ठेवून !!a{surjection} $f' \colon A \to \ran{f}$
मिळते. उघडच, ~$f$ हे !!{injective} असेल, तर~$f'$ हे !!{bijective}
असते. म्हणून \olref{prop:bijection-inverse} नुसार त्याला एकमेव
व्युत्क्रम असतो. संकेतलेखनात अगदी थोडी सवलत घेऊन आपण कधी कधी
$f'$ च्या व्युत्क्रमाला फक्त ``~$f$ चा व्युत्क्रम'' म्हणतो.
\end{explain}

\begin{prop}\ollabel{prop:left-right}%
  $f\colon A \to B$ ला डावा व्युत्क्रम~$g$ आणि उजवा व्युत्क्रम~$h$
  असल्यास $h = g$ असते, हे दाखवा.
\end{prop}

\begin{proof}
  सरावासाठी.
\end{proof}

\begin{prob}
  \olref[sfr][fun][inv]{prop:left-right} सिद्ध करा.
\end{prob}

\begin{prop}\ollabel{prop:inverse-unique}
प्रत्येक फलन~$f$ ला जास्तीत जास्त एक व्युत्क्रम असतो.
\end{prop}

\begin{proof}
  $g$ आणि $h$ हे दोन्ही~$f$ चे व्युत्क्रम आहेत असे समजा. मग विशेषतः
  ~$g$ हा~$f$ चा डावा व्युत्क्रम आणि~$h$ हा उजवा व्युत्क्रम असतो.
  \olref{prop:left-right} नुसार $g = h$.
\end{proof}

\end{document}
